\section{Related Work}
\label{sec:related}

\subsection{Transformers and Pretrained Language Models}

The Transformer architecture \citep{vaswani2017attention} replaced recurrent
sequence processing with attention over token states and remains the dominant
substrate for pretrained language models. BERT \citep{devlin2019bert} showed
the effectiveness of large-scale bidirectional pretraining for language
understanding, while GPT-style autoregressive language models demonstrated
few-shot behavior at scale \citep{brown2020language}. Parameter-efficient
adaptation methods such as LoRA \citep{hu2022lora} show that pretrained models
can often be modified without updating all backbone parameters.

PAT-ER builds on this pretrained-decoder substrate but does not treat the
decoder state as the only semantic carrier. Its central comparison is therefore
not ``pretrained versus not pretrained'', but ``token-pooled labels versus
side-state registers on the same pretrained backbone''.

\subsection{Semantic Roles and Event Structure}

Semantic-role theory studies how predicates structure their arguments. Dowty's
proto-role analysis \citep{dowty1991proto} argues that roles such as agent and
patient are better understood through clusters of entailments than through a
small rigid inventory. PropBank \citep{palmer2005propbank} and FrameNet
\citep{baker1998framenet} operationalized predicate-argument annotation for
computational NLP. These resources motivate the event-role side of PAT-ER: the
model should not only classify a primitive but also bind it to the event and
argument roles that license it.

\subsection{Logical Reasoning Benchmarks}

Natural-language reasoning datasets provide structured labels for entailment,
contradiction, and uncertainty. SNLI \citep{bowman2015snli} introduced a large
NLI benchmark with entailment, contradiction, and neutral labels. FEVER
\citep{thorne2018fever} framed fact verification as claim evaluation against
evidence. ProofWriter \citep{tafjord2021proofwriter} is closer to the present
work because it includes proofs, implications, and abductive statements over
templated natural language. FOLIO \citep{han2024folio} provides natural-language
premises with first-order logic annotations and verified reasoning labels.

PAT-ER uses these sources not to copy chain-of-thought into the input, but to
derive supervised side-state targets: primitive class, support state, rule
depth, entailment/refutation status, evidence pointers, and role-to-primitive
bridges. The no-chain-of-thought constraint is important: proof metadata is used
as supervision, not as rendered reasoning text.

\subsection{Hybrid Knowledge Extraction and Neural--Symbolic Reasoning}

Recent MDPI literature frames LLM reasoning as part of a wider effort to combine
statistical language models with structured knowledge. Dehal et al.
\citep{dehal2025kgllm} survey the reciprocal relationship between LLMs and
knowledge graphs, emphasizing that structured representations can improve
transparency, factual consistency, and reliability while LLMs automate entity,
relation, and schema extraction. Ivanisenko et al.
\citep{ivanisenko2024knowledgeextraction} demonstrate a domain-specific hybrid
pipeline in which ontology models, graph neural networks, and fine-tuned LLMs
are combined for scientific knowledge extraction. Liang et al.
\citep{liang2025neurosymbolic} review neural--symbolic reasoning architectures
and identify symbolic--continuous alignment, dynamic rule learning, and unified
architectures as open problems.

PAT-ER is aligned with this direction but takes a different design point. Rather
than attach a separate knowledge graph to a frozen text generator, it introduces
small event-role and primitive registers as first-class latent state inside the
decoder. The predicate-decomposition result of Tian and Meng
\citep{tian2024pdec} is especially relevant: it shows that predicate granularity
and predicate polysemy matter for interpretable knowledge-graph reasoning. PAT-ER
applies a related principle to language-model internals: logical primitives are
treated as a fixed predicate-like inventory whose semantic distinctions can be
learned, ablated, and measured.

\subsection{Structured Neural State and Tool Use}

Structured neural computation has a long history, from graph neural networks
\citep{scarselli2009gnn} to neural modules that impose explicit intermediate
structure. Tool-use work such as Toolformer \citep{schick2023toolformer} and
ReAct \citep{yao2023react} shows the value of coupling language models to
external actions. PAT-ER differs in two ways. First, the side-state is designed
around semantic primitive flow rather than around textual reasoning traces.
Second, the tool interface is decoupled from the architecture evaluation path:
base mode measures side-state, while \texttt{interface\_mode} handles
schema-grounded generation.
